\documentclass{article}
    \usepackage[english]{babel}
    \usepackage{geometry}
    \RequirePackage[nomessages]{fp}
    % Dimensioni derivate
    \newcommand{\dimensioni}{%
    % Dimensioni pagina
    \FPeval{PageHeightiii}{\PageHeight/3}
    \FPeval{PageHeightii}{\PageHeight-\PageHeight/3}
    \FPeval{PageWidthiii}{\PageWidth/3}
    \FPeval{PageWidthii}{\PageWidth-\PageWidth/3}
    \FPeval{PageWidthAP}{round(\PageWidth:2)}
    \FPeval{PageHeightAP}{round(\PageHeight:2)}
    % Rientro Lettrine
    \FPeval{Lettrine}{1/3}
    % Indent
    \FPeval{ParIndent}{\BaseLineSkip}
    \FPeval{ParIndentAP}{round(\ParIndent:2)}
    % Margini posizioni relative al bordo vicino
    \FPeval{LeftMargin}{round(\xLeftMargin:2)}
    \FPeval{RightMargin}{round(\PageWidth-\xRightMargin:2)}
    \FPeval{BottomMargin}{round(\yBottomMargin:2)}
    \FPeval{TopMargin}{round(\PageHeight-\yTopMargin:2)}
    % Altezza e larghezza del testo
    \FPeval{xyTextWidth}{\PageWidth-\LeftMargin-\RightMargin}
    \FPeval{xyTextHeight}{\PageHeight-\BottomMargin-\TopMargin}
    \FPeval{TextWidth}{round(\xyTextWidth:2)}
    \FPeval{TextHeight}{round(\xyTextHeight:2)}
    \FPeval{TextHeightiii}{\TextHeight/3}
    \FPeval{TextHeightii}{\TextHeight-\TextHeight/3}
    \FPeval{TextWidthiii}{\TextWidth/3}
    \FPeval{TextWidthii}{\TextWidth-\TextWidth/3}
    % Separazione della pagina
    \FPeval{yFootSkip}{\BottomMargin/(3/2)}
    \FPeval{FootSkip}{round(\BottomMargin-\yFootSkip:2)}
    % Dimensioni dei font
    \FPeval{FontSizePage}{\FontSize*2/3}
    \FPeval{FontSizeTitle}{\FontSize*5/3}
    \FPeval{FontSizeAuthor}{\FontSize*4/3}
    \FPeval{BaseLineTitle}{\FontSizeTitle*1.2}
    \FPeval{FontSizeAP}{round(\FontSize:2)}
    \FPeval{FontSizePageAP}{round(\FontSizePage:2)}
    \FPeval{FontSizeTitleAP}{round(\FontSizeTitle:2)}
    \FPeval{FontSizeAuthorAP}{round(\FontSizeAuthor:2)}
    \FPeval{BaseLineSkipAP}{round(\BaseLineSkip:2)}
    \FPeval{BaseLineTitleAP}{round(\BaseLineTitle:2)}
    % Titolo
    \FPeval{TitoloWidth}{\PageWidth-\LeftMargin*2}
    \FPeval{xTitolo}{\LeftMargin}
    \FPeval{yTitolo}{\PageHeight/9}
    \FPeval{xLogo}{\LeftMargin}
    \FPeval{yLogo}{\PageHeight*8/9-\TextWidth*0.16180/2}
    % Crediti
    \FPeval{xCrediti}{\PageWidthiii}
    \FPeval{CreditiWidth}{\PageWidthii-\LeftMargin}
    % Immagine
    \FPeval{yImmagine}{\PageHeight*3/9}
    % Separazione della virgola finale
    \FPeval{BaseLineSkipV}{5*\BaseLineSkip}
    }

    \newcommand{\dimePA}{%
    \FPeval{PageWidth}{210}%
    \FPeval{PageHeight}{297}%
    \FProot\solution{385}{2}%
    \FPeval{xLeftMargin}{1242150/14701-44100/14701*\solution}%
    \FPeval{xRightMargin}{1845060/14701+44100/14701*\solution}%
    \FPeval{yBottomMargin}{1756755/14701-10-62370/14701*\solution}%
    \FPeval{yTopMargin}{15+2609442/14701+62370/14701*\solution}%
    \FPeval{FontSize}{15}%
    \FPeval{BaseLineSkip}{18.35}%
    \dimensioni}

    \dimePA
    \geometry{paperwidth=\PageWidth mm,paperheight=\PageHeight mm,left=\LeftMargin mm,right=\RightMargin mm,top=\TopMargin mm,bottom=\BottomMargin mm,footskip=\FootSkip mm}
    \setlength{\parindent}{0pt}

    \usepackage{tikz}
    \usepackage{graphicx}
    \usepackage{calc}

    \usepackage[cmintegrals,cmbraces]{newtxmath}
    \usepackage{ebgaramond-maths}
    \usepackage[T1]{fontenc}
    \usepackage{anyfontsize}
    \usepackage{changepage}

    \newcommand{\titolo}[1]{\bgroup\rule{\linewidth}{1pt}\\[-14pt]
    \null\hfill\fontsize{10pt}{12pt}\sc\selectfont\strut#1\\[7pt]\egroup}

    \newcommand{\lavorofisio}[5]{%
    {\fontsize{14pt}{17pt}\selectfont\strut#1}\hfill\strut#2\\[-7pt]
    \begin{adjustwidth}{28pt}{28pt}
    \fontsize{10pt}{12pt}\selectfont\strut\textit{#3}\\[2pt]
    \fontsize{10pt}{12pt}\selectfont\strut#4\\[10pt]
    \end{adjustwidth}}


    \newcommand{\docenza}[5]{%
    {\fontsize{14pt}{17pt}\selectfont\strut#1}\hfill\strut#2\\[-7pt]
    \begin{adjustwidth}{28pt}{0pt}
    \fontsize{10pt}{12pt}\selectfont\strut\textit{#3}\\
    \end{adjustwidth}
    \vspace{-3.5pt}
    \begin{adjustwidth}{-21pt}{0pt}
    \begin{tabular}{@{}r@{\hspace{7pt}}p{\linewidth-56pt}}
    \fontsize{10pt}{12pt}\selectfont\strut Target: & #4\\
    \fontsize{10pt}{12pt}\selectfont\strut Contenuti: &#5\\[7pt]
    \end{tabular}
    \end{adjustwidth}
    }



    \newcommand{\paragrafo}[1]{{\fontsize{10pt}{12pt}\selectfont#1}\\[7pt]}

    \newcommand{\diploma}[4]{{\fontsize{14pt}{17pt}\selectfont#1}\hfill#2\\
    \begin{adjustwidth}{28pt}{28pt}
    \it #3\\
    \end{adjustwidth}}

    \newcommand{\universitaPD}[5]{#1\hfill#2\\#3\\#4\\ #5}

    \newcommand{\universitaMIB}[7]{{\fontsize{14pt}{17pt}\selectfont#1}\hfill#2\\
    \begin{adjustwidth}{28pt}{28pt}
    3: #3\\
    4: #4\\
    5: #5\\
    6: #6\\
    7: #7
    \end{adjustwidth}}


    \newcommand{\universita}[4]{{\fontsize{14pt}{17pt}\selectfont#1}\hfill#2\\
    \begin{adjustwidth}{28pt}{28pt}
    \textit{#3}\\
    #4\\[14pt]
    \end{adjustwidth}}


\begin{document}\raggedright\pagestyle{empty}
%\tikz[remember picture,overlay]{
%\draw (current page.north west) -- (current page.south east);
%\draw (current page.north east) -- (current page.south west);
%\draw ([yshift={0.5*\PageHeight mm-0.5*\PageWidth mm}]current page.center) circle (0.5*\PageWidth mm);
%\draw ([yshift={-0.5*\PageHeight mm+0.5*\PageWidth mm}]current page.center) circle (0.5*\PageWidth mm);
%}
\bgroup{\fontsize{20pt}{20pt}\selectfont
{\sc Elio A. Farina}\hfill \textit{R\'esum\'e}\\
\fontsize{10pt}{12pt}\sc\selectfont
44 Collinward Gardens\\
Newtownabbey, Co Antrim BT36 6DS\\
United Kingdom}\\
elio.farina@gmail.com\\
UK mobile: +447824705950\\
\egroup

\vfill
\titolo{Current occupations}

\lavorofisio{{CCO} --- Chief Content Officer}{From December 2015}{UC~Academy, Bree~29A2 3074 BB, Rotterdam, The Netherlands}{Editing and management of the \textit{Study guides} contents and all teaching material using \LaTeX\ and~{\sc Html} language.}{}

\vspace{20pt}
\lavorofisio{Digital typographer -- \LaTeX\ programmer}{From May 2014}{Freelancer for various clients all around the World}{\LaTeX\ language coding for publishing documents and report creation. Creation classes and packages for custom made document templates. Drawing vectorial images using TikZ/\LaTeX. Teaching \LaTeX\ and basics of \texttt{git}. Assisting clients on \LaTeX\ code related issues and debugging.\\
Portfolio at \texttt{http://badroomtales.me/portfolio/EAFportfolio.pdf}.}{}

\vfill
\titolo{Previous occupation}
\lavorofisio{Physiotherapist}{From December 2009 to April 2014}{Freelancer.}{Physiotherapy treatments for various typology of patients, private and for the Italian public health system.}

\vfill
\titolo{School}
%\paragrafo{Scuola superiore}
%\diploma{Diploma di maturit\`a scientifica}{1999}{Liceo Scientifico Statale {\em ``Paolo Frisi''} di Monza.}{diploma di maturit\`a scientifica.}

%\vspace{21pt}
%\paragrafo{Universit\`a}

%\universita{Astronomia}{Da ottobre 1999 ad agosto 2004}{Universit\`a degli Studi di Padova, Facolt\`a di Scienze {\sc mm. ff. nn.}, Dipartimento di Astronomia,\\ {Corso di Laurea in Astronomia} Vecchio Ordinamento.}{Ho frequentato tutti i corsi dei quattro anni del Vecchio Ordinamento del Corso di Laurea in Astronomia e superato con profitto gli esami del primo biennio.}


\universita{Physiotherapist}{From October 2006 to November 2009}%
{Universit\`a degli Studi di Milano~--~Bicocca, Facolt\`a di Medicina,\\ Corso di Laurea in Fisioterapia abilitante alla professione di Fisioterapista.}%
{Degree in Physiotherapy, 11 November 2009.}

\vfill
\null\hfill Newtownabbey, \today\\[-7pt]
\null\hfill\begin{minipage}{0.5\linewidth}\centering
\resizebox{0.8\linewidth}{!}{\begin{tikzpicture}[y=1pt,x=1pt,yscale=-1,line join=miter,line cap=round,miter limit=4.00,line width=4pt]
% Eli
  \path[draw=black]
    (210,21) .. controls
    (194,7) and (167,34) .. (167,49) .. controls
    (166,75) and (233,86) .. (251,60) .. controls
    (261,47) and (215,53) .. (201,65) .. controls
    (179,83) and (185,97) .. (212,91) .. controls
    (296,70) and (43,76) .. (20,128) .. controls
    (-3,182) and (334,77) .. (298,48) .. controls
    (234,-3) and (73,90) .. (79,140) .. controls
    (84,182) and (194,157) .. (228,148) .. controls
    (245,144) and (252,157) .. (270,160) .. controls
    (282,162) and (281,146) .. (281,146);
% o
  \path[draw=black] (329,146) .. controls (282,63) and (260,192) .. (350,123);
% A
  \path[draw=black]
    (431.3630,160.5543) .. controls
    (461.8548,122.9296) and (494.5664,72.4338) .. (512.3767,16.5842) .. controls
    (500.6766,70.3928) and (489.1791,161.0044) .. (517.2325,197.8225) .. controls
    (501.8808,92.9720) and (429.2264,95.7237) .. (383.3826,103.6564) .. controls
    (472.4962,86.2726) and (503.7766,77.7906) .. (560.7513,77.7906);
% F
  \path[draw=black]
    (698.8650,36.7957) .. controls
    (676.3174,71.4507) and (674.2967,108.7374) .. (678.8555,148.5556) .. controls
    (748.2250,69.0229) and (611.8115,-20.2442) .. (561.7274,44.1163) .. controls
    (534.2121,79.4746) and (746.3889,75.4141) .. (810.1368,40.2120) .. controls
    (728.5435,45.8691) and (583.5430,95.0472) .. (583.6889,109.0248) .. controls
    (583.9058,129.8135) and (714.9701,104.8998) .. (768.1658,94.8718);
% arin
  \path[draw=black]
    (741.9484,159.5460) .. controls
    (613.3203,183.1078) and (714.1827,95.4374) .. (750.3144,159.6788) .. controls
    (821.3610,134.1986) and (885.1179,144.1260) .. (914.0065,152.7279);
% a
  \path[draw=black]
    (962.8914,157.8282) .. controls
    (954.7155,93.0862) and (864.7374,162.4792) .. (949.7145,163.6846) .. controls
    (1011.1273,164.5558) and (987.5983,120.9472) .. (1125.8949,165.1487);
% l.i
  \fill (285,101) circle (4pt);
% A.
  \fill (541,163) circle (4pt);
% r.i
  \fill (872,87) circle (4pt);
\end{tikzpicture}
}\\[-16pt]
\rule{\linewidth}{0.5pt}
\end{minipage}

\end{document}
